\documentclass[10pt,a4paper]{amsart}
\usepackage[margin=19mm]{geometry}
\usepackage{amsmath,amssymb,mathtools}
\usepackage[scr=rsfs,cal=euler]{mathalfa}
\usepackage{xfrac}
\usepackage[unicode,hidelinks,bookmarks=false]{hyperref}
\newcommand{\ie}{i.e.\ }
\newcommand{\cf}{cf.\ }
\newcommand{\resp}{respectively\ }
\newcommand{\Subsection}{\subsection}
\providecommand{\nobreakdash}{\nobreak\hbox{-}}
\setlength{\emergencystretch}{2em}
\multlinegap0pt
\usepackage{amssymb}
\usepackage{xcolor}
\newtheorem{theorem}{Theorem}
\title{Lower order term for the fractional Laplacian with a Hardy potential}
\author{Rub\'en Fi\~nana, Alexis Molino}
\date{Source: arXiv:2608.26828v1 (CC BY 4.0)}
\begin{document}
\maketitle

\noindent\emph{Source and licence.} Lower order term for the fractional Laplacian with a Hardy potential. Version arXiv:2608.26828v1 (CC BY 4.0), distributed by arXiv under the Creative Commons Attribution 4.0 International licence, http://creativecommons.org/licenses/by/4.0/, as recorded in the arXiv licence metadata for this exact version and shown on the abstract page https://arxiv.org/abs/2608.26828v1. Full licence notice: \texttt{licenses/arxiv-2608.26828-CC-BY-4.0.txt}.\par\smallskip

\noindent\textit{Editorial scope.} Counted unit: the article's theorem teo1 (source0.tex lines 295--298). The article's complete original proof of it is delivered with it (source0.tex lines 299--364, one \texttt{proof} environment of 66 lines). No mathematics is written, repaired or re-derived: the statement and the proof are the article's own lines, in the article's own notation, and no retained line is substituted or re-flowed.  Retained uncounted as the source-local context the proof needs: lines 188--190; lines 193--198; lines 200--202; lines 245--248; lines 274--276; lines 291--293.  Nothing was omitted from any retained span, so the package carries no omission record.  No retained line contains a citation, so the wrapper carries no bibliography and no bibliography entry.  No retained line contains a figure environment, an image-inclusion command or drawing code, so no image is delivered and the package declares no asset dependency.  Editorial formatting only, in addition: an AMS \texttt{amsart} preamble that restates the article's own declarations and macros used by the retained lines, namely the environments \texttt{theorem} on separate counters, together with the standard packages those lines need.  The retained environments print in this document as Theorem 1 in order of appearance, whereas the article prints the same items with the numbering of its own full document.  The complete native source of the article as distributed in the arXiv e-print for this exact version is bundled unchanged as \texttt{sources/208666/source0.tex}.
		\begin{equation}\label{DesigualdadHardy}
			\frac{c_{N,s}}{2}\int_{\mathbb{R}^N}\int_{\mathbb{R}^N} \frac{|u(x)-u(y)|^2}{|x-y|^{N+2s}}\geq \Lambda_{N,s}\int_{\Omega} \frac{u^2}{|x|^{2s}}
		\end{equation}		

		\begin{equation}\label{problemaoriginal}
			\begin{cases}
				\displaystyle(-\Delta)^s u+ g(x)|u|^{p-1}u= \lambda \frac{u}{|x|^{2s}}+f(x), &  \text{in $\Omega$,}\\
				\mathnormal{u}=0, & \text{in $\mathbb{R}^N\setminus\Omega$,}
			\end{cases}
		\end{equation}

		\begin{equation}\label{condicion}
			\frac{1}{|x|^{2s}g^{\frac{2}{p+1}}}\in L^{\frac{p+1}{p-1}}(\Omega).
		\end{equation}

		\begin{equation}\label{DualE}
			\langle v_f, h
			\rangle = \int_{\Omega} fhg,\textit{ $\forall h\in L^{p+1}_g(\Omega)$.}
		\end{equation}

		\begin{equation}\label{acotacion}
			\int_\Omega \frac{u^2}{|x|^{2s}} \leq C_1\left(\int_{\Omega}|u|^{p+1}g\right)^{\frac{2}{p+1}}<+\infty,
		\end{equation}

		\begin{theorem}\label{ekeland}
			If $X$ is a Banach space and $\phi\in C^1(X,\mathbb{R})$ is bounded from below, then there exists a minimizing sequence $u_n$ for $\phi$ such that $\phi(u_n)\rightarrow \inf_X(\phi)$ and $\phi'(u_n)\rightarrow 0$ in $X^*$ as $n\rightarrow+\infty$.
		\end{theorem}

		\begin{theorem}\label{teo1}
			Let $\Omega\subset \mathbb{R}^N$ be a smooth bounded domain such that $0\in \Omega$,  $f \in L^{(p+1)/p}_g(\Omega)$ with $g$ satisfying  inequality \eqref{condicion}. Then, for every $\lambda\in \mathbb{R}$ there exists a solution $u\in E$ to the problem \eqref{problemaoriginal}. 
			
		\end{theorem}

		\begin{proof}
			With the aim of guaranteeing the existence of a solution,  the coercivity and the lower boundedness of the functional $\phi_{\lambda}$ are needed.\par
			There are two possible scenarios depending on the value of $\lambda$:
			\begin{enumerate}
				\item[a)] Case $\lambda>0$.\par 
				Because of \eqref{acotacion} and \eqref{DualE} for $u\in E$, 
				\begin{align*}
					\phi_{\lambda} (u)\geq&  \frac{1}{2}\|u\|^2_{H^s_0(\Omega)}+\frac{1}{p+1}\int_{\Omega}|u|^{p+1}g-\frac{\lambda}{2}C_1\left(\int_{\Omega}|u|^{p+1}g\right)^{\frac{2}{p+1}}\\&-\langle f,u\rangle\\
					\geq& \frac{1}{2}\|u\|^2_{H^s_0(\Omega)}+\frac{1}{p+1}\int_{\Omega}|u|^{p+1}g-\frac{\lambda}{2} C_1\left(\int_{\Omega}|u|^{p+1}g\right)^{\frac{2}{p+1}}\\&-\| f\|_{E^*}\|u\|_{E}\\
					=&\frac{1}{2}\|u\|^2_{H^s_0(\Omega)}+\frac{1}{p+1}|u|^{p+1}_{L^{p+1}_g(\Omega)}-\frac{\lambda}{2}C_1|u|^2_{L^{p+1}_g(\Omega)}-\| f\|_{E^*}\|u\|_{E}.
				\end{align*}
				
				\item[b)] Case $\lambda\leq 0$.\par 
				It is enough to use \eqref{DualE} for $u\in E$, 
				\begin{align*}
					\phi_{\lambda} (u)&= \frac{1}{2}\|u\|^2_{H^s_0(\Omega)}+\frac{1}{p+1}\int_{\Omega}|u|^{p+1}g-\frac{\lambda}{2} \int_{\Omega}\frac{u^2}{|x|^{2s}}-\langle f,u\rangle\\
					&\geq \frac{1}{2}\|u\|^2_{H^s_0(\Omega)}+\frac{1}{p+1}|u|^{p+1}_{L^{p+1}_g(\Omega)}-\| f\|_{E^*}\|u\|_{E}.
				\end{align*}
			\end{enumerate}
			So for $\lambda\in \mathbb{R}$, $\phi_{\lambda}$ is coercive, yielding its lower boundedness. The claim is thus proved.
			
			
			Consequently, Theorem \ref{ekeland} can be applied to obtain a minimizing sequence, denoted by $u_n\in E$, such that 
			\[
			\phi_{\lambda}(u_n)\rightarrow \inf_{E}{\phi_{\lambda}},
			\]
			and for all $w\in E$, \begin{equation}\label{convergencia}
				\langle \phi'_{\lambda}(u_n),w\rangle\rightarrow 0,
			\end{equation}
			where 
			\begin{align*}
				\langle \phi'_{\lambda}(u_n),w\rangle=&\int_{\mathbb{R}^N}(-\Delta)^{\frac{s}{2}}u_n(-\Delta)^{\frac{s}{2}}w + \int_{\Omega}|u_n|^{p-1}u_nwg
				\\&-\lambda \int_{\Omega}\frac{u_n}{|x|^{2s}}w-\int_{\Omega}fwg.
			\end{align*}
			Due to the coercivity of the functional, $u_n$ is bounded in $E$, thus there exists a subsequence, still denoted by $u_n$, such that $u_n\rightharpoonup u$ in $E$. Therefore,
			\begingroup 
			\renewcommand{\theenumi}{(\alph{enumi})}
			\renewcommand{\labelenumi}{\theenumi}
			\begin{enumerate}
				\item $u_n\rightharpoonup u$ in $H^s_0(\Omega)$\label{conv1}.
				\item $u_n\rightharpoonup u$ in $L^{p+1}_g(\Omega)$\label{conv2}.
			\end{enumerate}
			\endgroup
			Taking into account \eqref{convergencia}, the limit $u$ is a solution of \eqref{problemaoriginal} if the following claim is satisfied
			\begin{equation}\label{limitsolution}
				\langle \phi'_{\lambda}(u_n),w\rangle\rightarrow \langle \phi'_{\lambda}(u),w\rangle= 0.
			\end{equation}
			Indeed, the convergence \ref{conv1} directly implies that 
			\begin{equation}\label{limitsolution1}
				\int_{\mathbb{R}^N}u_n(-\Delta)^{s}w\rightarrow \int_{\mathbb{R}^N}u(-\Delta)^{s}w.
			\end{equation}
			Moreover, by \ref{conv1}, $u_n\rightarrow u$ in $L^m(\Omega)$ for $m\in [1,2^*_s)$, so $u_n \rightarrow u$ a.e., therefore  $|u_n|^{p-1}u_n\rightarrow |u|^{p-1}u$ a.e.
			On the other hand, due to \ref{conv2} $u_n$ is bounded in $L^{p+1}_g(\Omega)$, thus $|u_n|^{p-1}u_n g^{p/(p+1)}$ is bounded in $ L^{(p+1)/p}(\Omega)$. That fact, along with convergence a.e., implies that $|u_n|^{p-1}u_n g^{p/(p+1)}\rightharpoonup |u|^{p-1}u g^{p/(p+1)}$ in $L^{(p+1)/p}(\Omega)$. As a consequence
			\begin{equation}\label{limitsolution2}
				\int_{\Omega}|u_n|^{p-1}u_nwg\rightarrow  \int_{\Omega}|u|^{p-1}uwg.
			\end{equation}
			Finally, consider $G_w: H^s_0(\Omega)\rightarrow \mathbb{R}$ defined as
			\[G_w(u)=\int_{\Omega}\frac{u}{|x|^{2s}}w.\]
			Notice that, if the H\"older inequality and Hardy inequality \eqref{DesigualdadHardy} are applied,
			\[\int_{\Omega}\frac{u}{|x|^{2s}}w\leq \left\|\frac{u}{|x|^{s}}\right\|_{L^2(\Omega)} \left\|\frac{w}{|x|^{s}}\right\|_{L^2(\Omega)} \]\[\leq \frac{1}{\Lambda_{N,s}}\|u\|_{H^s_0(\Omega)}\|w\|_{H^s_0(\Omega)}.\]
			Hence, $G_w$ is a continuous linear functional and, by \ref{conv1}, ${G_w(u_n)\rightarrow G_w(u)}$. This implies that
			\begin{equation}\label{limitsolution3}
				\int_\Omega \frac{u_n}{|x|^{2s}}w\rightarrow \int_\Omega \frac{u}{|x|^{2s}}w.
			\end{equation}
			As a result of the limits \eqref{limitsolution1},  \eqref{limitsolution2} and  \eqref{limitsolution3}, the claim \eqref{limitsolution} is proved.
		\end{proof}

\end{document}
